\documentclass{article}
\usepackage[T1]{fontenc}
\usepackage{lmodern}
\usepackage{graphicx}
\usepackage{amsmath,amssymb}
\usepackage[a4paper,margin=2cm]{geometry}
\usepackage[absolute,overlay]{textpos}
\setlength{\TPHorizModule}{1mm}
\setlength{\TPVertModule}{1mm}
\usepackage[indonesian]{babel}
\usepackage{hyperref}
\setlength{\emergencystretch}{2em}
% unit: Halaman 13

\begin{document}

% === Halaman 13 (python/native) ===

13

• Secara umum, jika pergeseran huruf sejauh k, maka:


 Enkripsi: c = E(p) = (p + k) mod 26

 Dekripsi: p = D(c) = (c – k) mod 26


       k = kunci rahasia

• Untuk alfabet berupa 256 karakter ASCII, maka:


 Enkripsi: c = E(p) = (p + k) mod 256

 Dekripsi: p = D(ci) = (c – k) mod 256

                   k = kunci rahasia

\end{document}
